\documentclass[11pt]{article}
\usepackage{amsmath,amssymb,amsthm,hyperref}
\newtheorem{theorem}{Theorem}
\newtheorem{lemma}[theorem]{Lemma}
\title{First-Order Rigidity of a Boolean Permutation Construction}
\author{}
\date{November 23, 2026}
\begin{document}
\maketitle

\paragraph{Scope.}
This note audits a specific possible use of the Boolean-only moment
construction in \texttt{ideal-moment-rank-audit.tex}. It does not prove
the requested universal head bound. The results below have elementary
manual proofs; they have not been formalized in Lean.

\section{The permutation forms}

For $x\in\{0,1\}^n$, define
\[
 u_0(x)=\sum_{i=1}^n x_i,\qquad
 u_k(x)=k(1-x_k)+\sum_{i=k+1}^n x_i\quad(1\le k\le n).
\]
Their values are a permutation of $0,1,\ldots,n$ at every vertex.
Induct on $n$: the last entry is $n$ when $x_n=0$, and the earlier
entries are the smaller construction. When $x_n=1$, the last entry
is zero and each earlier entry is increased by one.

\begin{lemma}[Affine rigidity on the level sets]
Suppose $n\ge4$ and $r:\{0,1\}^n\to\mathbb R$ has the following
property: for every $k,j\in\{0,\ldots,n\}$, the restriction of $r$
to $\{x:u_k(x)=j\}$ agrees with an affine function of $x$.
Then $r$ agrees with a single affine function on the entire cube.
\end{lemma}
\begin{proof}
Because $u_n=n(1-x_n)$, the restrictions of $r$ to both faces
$x_n=0$ and $x_n=1$ are affine. Subtracting these two affine
expressions shows that, on the cube,
\[
 r(x)=L(x)+x_n\sum_{i=1}^{n-1}c_i x_i
\]
for some affine function $L$ and real numbers $c_i$.

The level set $u_0=2$ is the Hamming layer of weight two. Consequently
$x_n\sum_{i<n}c_i x_i$ agrees on this layer with an affine expression
$a+\sum_{i=1}^n w_i x_i$. For all distinct $i,j<n$, evaluating at
$e_i+e_j$ gives $a+w_i+w_j=0$. Since $n-1\ge3$, comparing these
equalities shows that all $w_i$, $i<n$, are equal to some $w$ and
$a=-2w$. Evaluation at $e_i+e_n$ then gives
$c_i=a+w_i+w_n=w_n-w$. Thus every $c_i$ equals one common value $c$.

Fix $i\in\{1,\ldots,n-3\}$ and set $k=n-2$. For $t\in\{0,1\}$
consider the four vertices
\[
 p_t=t e_i+e_k+e_{n-1},\qquad
 q_t=t e_i+e_k+e_n.
\]
At all four vertices $u_k=1$. Hence the level-set hypothesis implies
\[
 r(q_1)-r(q_0)=r(p_1)-r(p_0),
\]
because the same affine expression represents $r$ there and both
differences are taken in direction $e_i$.
The affine term $L$ has equal differences as well. But the corresponding
differences of $c x_n\sum_{j<n}x_j$ are $c$ on the $q$ pair and zero
on the $p$ pair. Therefore $c=0$, as required.
\end{proof}

\section{Application to product differences}

Let $a_0,\ldots,a_n,b_0,\ldots,b_n$ be arbitrary real affine forms.
Introduce a deformation parameter $\varepsilon$ and an auxiliary
variable $s$:
\[
 F_\varepsilon(s,x)=\prod_{k=0}^n(s+u_k(x)+\varepsilon a_k(x)),
 \qquad
 G_\varepsilon(s,x)=\prod_{j=0}^n(s+j+\varepsilon b_j(x)).
\]
At $\varepsilon=0$ these two products agree on the cube for every $s$.

\begin{theorem}[First-order ideal moment differences are affine]
Suppose $n\ge4$ and
\[
 \left.\frac{\partial}{\partial\varepsilon}
 (F_\varepsilon(s,x)-G_\varepsilon(s,x))\right|_{\varepsilon=0}
 =r(x)
\]
for every Boolean vertex $x$ and every real $s$, where the right side
is independent of $s$. Then $r$ agrees with an affine form on the cube.
Conversely, every affine form $r$ occurs as such a first-order difference.
\end{theorem}
\begin{proof}
Put $P(s)=\prod_{j=0}^n(s+j)$ and
\[
 d_j=P'(-j)=(-1)^j j!(n-j)!,
\]
which is nonzero. Fix a Boolean vertex $x$. If $u_k(x)=j$, then
evaluation of the first variation at $s=-j$ leaves exactly one term
from each product, because the roots are distinct. It yields
\[
 r(x)=d_j\bigl(a_k(x)-b_j(x)\bigr).
\]
Thus $r$ agrees with an affine form on each level set $u_k=j$.
The lemma proves the first assertion.

For the converse, set $a_k=0$ for every $k$ and
$b_j=-r/d_j$ for every $j$. These are affine forms. The first
variation is
\[
 r(x)\sum_{j=0}^n\frac{1}{d_j}\prod_{\ell\ne j}(s+\ell)=r(x).
\]
Indeed the polynomial in $s$ given by the sum has degree at most $n$
and takes the value one at all $n+1$ distinct points $-j$. Subtract
one and apply the elementary root bound for polynomials.
\end{proof}

\paragraph{Precise limitation.}
The conclusion concerns the first variation at these particular two
factor lists. It is not asserted for arbitrary affine permutation cubes,
for nonconstant unperturbed product differences, or for higher-order
deformations. In particular the proof does not exclude a nonlinear
sign construction from higher-order cancellations. Such a construction
would require a separate argument controlling those cancellations and
their head count.

\section{All zero first variations}

The following classification will distinguish a genuine second-order
obstruction from a choice of an inconvenient first-order representative.
Write $p=x_{n-1}$ and $q=x_n$.

\begin{lemma}[Zero-variation normal form]
Let $n\ge3$. The first variation above is identically zero on the cube
for every $s$ if and only if, after subtracting one common affine form
$C(x)$ and a common multiple $\lambda u_k(x)$ from each $a_k$, and
$C(x)+\lambda j$ from each $b_j$, the remaining forms are
\[
 a_k=0\ (0\le k\le n-2),\quad
 a_{n-1}=\theta(p-q),\quad
 a_n=\theta p+(\theta+\beta+\delta)q-\theta-\delta,
\]
\[
 b_j=0\ (1\le j\le n-1),\quad
 b_0=\theta p+\beta q,\quad
 b_n=\theta(p-1)+\delta(q-1)
\]
for real constants $\theta,\beta,\delta$.
Every such zero first variation integrates to the exact identity
$F_\varepsilon=G_\varepsilon$ on the cube for every $s,\varepsilon$.
\end{lemma}
\begin{proof}
Evaluation at the roots, as before, shows that zero first variation is
equivalent to $a_k(x)=b_j(x)$ whenever $u_k(x)=j$.
Put $t=\sum_i x_i$. For $1\le j\le n-1$, the layer $t=j$ affinely
spans its defining hyperplane. One proof is to use differences of
weight-$j$ vertices differing by exchanging two coordinates: these
differences span $\{v:\sum_i v_i=0\}$.
The same statement holds after complementing coordinate one, for
the layer $u_1=t+1-2x_1=j$. Hence
\[
 a_0-b_j=\alpha_j(t-j),\qquad
 a_1-b_j=\eta_j(t+1-2x_1-j).
\]
Comparing coefficients of $x_1$ and any $x_i$ with $i>1$ shows that
$\alpha_j$ is independent of $j$, and so is $\eta_j$.
Comparing constants for the two distinct indices $j=1,2$ gives
$\alpha_j=\eta_j=\lambda$. Therefore
$C=a_0-\lambda t$ gives $a_0=C+\lambda u_0$,
$a_1=C+\lambda u_1$ and $b_j=C+\lambda j$ for all interior $j$.
Subtract these common contributions.

For $2\le k\le n-2$, the set where $1\le u_k\le n-1$ has
$2^n-2^k>2^{n-1}$ vertices, and the normalized $a_k$ vanishes there.
A nonzero affine function has at most $2^{n-1}$ Boolean zeros:
pair vertices in a coordinate with a nonzero coefficient; at most one
endpoint of each pair can be a zero. A nonzero constant has no zeros.
Thus these $a_k$ vanish identically as well.

On the face $p=q=0$, the row with value zero has index at most $n-2$,
so $b_0$ vanishes there. On the face $p=q=1$, the row with value $n$
has index at most $n-2$, so $b_n$ vanishes there. It follows that
$b_0=\alpha p+\beta q$ and
$b_n=\gamma(p-1)+\delta(q-1)$.
The form $a_{n-1}$ vanishes on both faces $p=q=0$ and $p=q=1$,
since the corresponding root values $n-1$ and $1$ are interior.
It is therefore $\theta(p-q)$. Matching the roots at $(p,q)=(1,0)$
and $(0,1)$ gives $\alpha=\gamma=\theta$.
Finally the row $u_n$ equals $n$ on $q=0$ and zero on $q=1$.
Its two restrictions give the displayed formula for $a_n$.

Conversely, substitution at the four choices of $(p,q)$ verifies
$a_k=b_{u_k}$ for the last two rows. For an earlier row with value
zero, its tail is all zero and hence $b_0=0$; for an earlier row with
value $n$, its tail is all one and hence $b_n=0$.
All other matches are zero after normalization. Restoring $C$ and
$\lambda$ preserves every match. Therefore the two perturbed lists
are still pointwise permutations of one another for every
$\varepsilon$, proving the exact identity and the converse.
\end{proof}

\section{A second-order sparsity condition}

\begin{theorem}[At most one active earlier coordinate]
Suppose $n\ge4$ and replace the factors in the preceding construction
by
\[
 u_k+\varepsilon a_k+\varepsilon^2 c_k,
 \qquad j+\varepsilon b_j+\varepsilon^2 e_j,
\]
where all the displayed coefficient functions are affine in $x$.
Assume that the coefficients of both $\varepsilon$ and
$\varepsilon^2$ in $F_\varepsilon-G_\varepsilon$ are independent of
$s$ on the Boolean cube. Write the first coefficient as
$r(x)=r_0+\sum_i r_i x_i$, using the first-order theorem.
Then
\[
 r_i r_\ell=0\qquad(1\le i<\ell\le n-2).
\]
In particular, this first coefficient depends on at most three input
coordinates, including the last two coordinates.
\end{theorem}
\begin{proof}
Adding $r/d_j$ to each $b_j$ makes the first variation zero.
Apply the normal-form lemma to those modified forms.
The common affine shift $C$ and scalar scaling $\lambda$ may be
removed: subtract $\varepsilon C(x)$ from every root and divide every
root by $1+\varepsilon\lambda$. This preserves affineness of each
coefficient function and independence of $s$ to the orders under
consideration. Indeed the new product difference is the old one
evaluated at $(1+\varepsilon\lambda)s-\varepsilon C(x)$ and divided
by $(1+\varepsilon\lambda)^{n+1}$.
Its first coefficient remains $r$; its second coefficient changes by
$-(n+1)\lambda r$.
Only formal series through order two are needed here.

We can consequently assume that $a_k=0$ for $k\le n-2$, both
remaining $a_k$ depend only on $(p,q)$, and
\[
 b_j=t_j(p,q)-r(x)/d_j,
\]
where every $t_j$ depends only on $(p,q)$ and $t_j=0$ for interior
$j$. Denote the second coefficient of the product difference by
$r^{(2)}(x)$.

At $s=-j$, the contributions from the second-order root coefficients
are $d_j(c_{u^{-1}(j)}-e_j)$. Let $Q_j(x)$ be the remaining
contribution, formed by pairs of first-order root coefficients.
It has the form
\[
 Q_j(x)=T_{j,0}(p,q)+T_{j,1}(p,q)r(x)+\kappa_j r(x)^2,
 \qquad
 \kappa_j=-\sum_{\ell\ne j}\frac1{d_\ell(\ell-j)}.
\]
For the first product only the two last rows can contribute, and
their root values and first coefficients depend only on $(p,q)$.
For the second product the term indexed by $j,\ell$ at $s=-j$ is
$d_j b_jb_\ell/(\ell-j)$. Expanding the two factors gives exactly
the claimed constant, linear, and quadratic terms in $r$, with the
minus sign from subtracting the second product.

The elementary partial-fraction identity for $1/P(s)$ gives
\[
 \kappa_j=\frac1{d_j}\sum_{\ell\ne j}\frac1{\ell-j}
          =\frac{H_{n-j}-H_j}{d_j},
 \qquad H_a=\sum_{v=1}^a\frac1v,\quad H_0=0.
\]
For detail, write $P(s)=(s+j)P_j(s)$ and subtract
$1/(d_j(s+j))$ from $1/P(s)$. Its value at $s=-j$ after cancellation
is $-P_j'(-j)/d_j^2$. Compare this with the sum of the other
partial fractions, and use
$P_j'(-j)/d_j=\sum_{\ell\ne j}1/(\ell-j)$.
In particular
\[
 \kappa_0-\kappa_1
 =\frac{(n+1)H_{n-1}+1/n-n}{n!}>0,
\]
since $H_{n-1}\ge1$.

On the face $p=q=1$, the rows $n$ and $n-1$ have root values zero
and one, respectively. The two root evaluations of $r^{(2)}$ give
\[
 r^{(2)}=d_0(c_n-e_0)+Q_0
        =d_1(c_{n-1}-e_1)+Q_1.
\]
Hence $Q_0-Q_1$ is affine in the first $n-2$ variables on this face.
The displayed form of $Q_j$ shows that
$(\kappa_0-\kappa_1)r(x)^2$ is also affine there.
Its mixed Boolean difference in two distinct earlier coordinates
$i,\ell$ is
$2(\kappa_0-\kappa_1)r_i r_\ell$. It must vanish.
The strict inequality for the scalar factor proves the assertion.
\end{proof}

\paragraph{Scope of the new obstruction.}
This is a necessary condition for a second-order continuation of the
specified ideal moment identity. It is not a classification of all
higher-order branches, all permutation constructions, or ordinary
head scores. A general affine first coefficient, although allowed
by the linearized equations, need not extend even to second order.

\section{Target-dependent level sets at arbitrary sizes}

There is a different obstruction which does not require the special
root list $0,\ldots,n$, but does require a constant reference list.

\begin{theorem}[Counting affine classifiers on projection level sets]
Let $N=2^n$, $n\ge1$ and $M\ge1$. Consider all strict Boolean sign
labelings that are affine on each level set of some affine projection
$L:\{0,1\}^n\to\mathbb R$ taking at most $M$ distinct values.
Here the projection and the affine classifier on each of its level
sets may all depend on the target labeling. The number of such
labelings is at most
\[
 N^{2n}\left[e\left(1+\frac{N}{M(n+1)}\right)\right]^{M(n+1)}.
\]
For fixed $C>0$, if $M\le C2^n/n^2$, this is $2^{o(2^n)}$ as
$n\to\infty$.
\end{theorem}
\begin{proof}
The partition into level sets is determined by which of the at most
$T=N(N-1)/2$ difference vectors $x-y$ are orthogonal to the slope
vector of $L$. The span of those orthogonal difference vectors has a
basis of at most $n$ members. Choose a canonical such basis from a
fixed ordering of all difference vectors, and pad it to length $n$
with a separate zero symbol. This encodes the full set of orthogonal
differences: it consists exactly of the listed difference vectors in
that basis span. Thus there are at most $(T+1)^n\le N^{2n}$ possible
level-set partitions. This includes the constant projection.

On a set of $m$ points, affine strict sign tests have at most
$\sum_{k=0}^{d}\binom{m}{k}$ sign patterns, where $d=n+1$.
To prove this elementary bound, take the $m$ hyperplanes in the
$d$-dimensional parameter space given by zero evaluations.
Every nonempty strict sign region is convex, so is one chamber.
Adding the last hyperplane splits at most as many chambers as the
previous hyperplanes cut on that hyperplane, a space of dimension
$d-1$. Induction gives the binomial-sum bound, including degeneracies.
When $m\ge d$, put $t=d/m\le1$ in the binomial theorem to obtain
\[
 \sum_{k=0}^{d}\binom{m}{k}
 \le t^{-d}(1+t)^m\le (em/d)^d.
\]
For $m<d$, the bound $2^m\le e^d$ is sufficient. In either case
the number of sign patterns is at most $[e(1+m/d)]^d$.

For a partition with $q\le M$ nonempty cells of sizes $m_1,\ldots,m_q$,
multiply these bounds and use concavity of $\log(1+t)$ to obtain
\[
 \prod_{j=1}^q[e(1+m_j/d)]^d
 \le[e(1+N/(qd))]^{qd}
 \le[e(1+N/(Md))]^{Md}.
\]
The last inequality follows because $z\log(e(1+N/z))$ is increasing
for $z>0$: its derivative is
$1+\log(1+N/z)-N/(N+z)>0$.
Multiplying by the number of partitions proves the explicit bound.

Its base-two logarithm is at most
$2n^2+M(n+1)\log_2(e(1+N/(M(n+1))))$.
The same monotonicity, with $M(n+1)\le2CN/n$, bounds this by
$2n^2+O_C(N\log n/n)=o(N)$.
\end{proof}

\begin{theorem}[Constant-reference permutation perturbations do not suffice]
Let $\lambda_1,\ldots,\lambda_M$ be distinct real numbers, and suppose
the affine forms $u_1,\ldots,u_M$ have this value multiset at every
Boolean vertex. Suppose affine perturbations $a_h,b_j$ satisfy
\[
 \left.\partial_\varepsilon
 \left(\prod_{h=1}^M(s+u_h+\varepsilon a_h)
 -\prod_{j=1}^M(s+\lambda_j+\varepsilon b_j)\right)
 \right|_{\varepsilon=0}=r(x)
\]
on the cube, independently of $s$. If $r$ has no Boolean zero, its
sign labeling belongs to the family in the preceding theorem.
Consequently, such first-order constructions with
$M\le C2^n/n^2$ do not represent every Boolean function, even when
both the root list and the permutation forms depend on the target.
\end{theorem}
\begin{proof}
Fix one row $h$. On its level set $u_h=\lambda_j$, evaluation at
$s=-\lambda_j$ gives
\[
 r(x)=\left(\prod_{\ell\ne j}(\lambda_\ell-\lambda_j)\right)
       (a_h(x)-b_j(x)).
\]
The prefactor is nonzero and constant on this level set, so the
strict signs agree there with an affine classifier. The single
projection $u_h$ takes at most $M$ values. Apply the counting theorem.
\end{proof}

\paragraph{A distinction essential to this counting argument.}
The second factor list at the base point is constant. If it is another
nonconstant affine permutation list, the index of its row carrying
$\lambda_j$ can vary within the level set of the selected first row.
Then the displayed affine-on-level-sets conclusion does not follow.
The count is not asserted for that case, for higher-order leading
differences, or for unrestricted sums of affine fractions.

\section{Diagnostic and audit}

The script \texttt{probe\_boolean\_moment\_tangent.py} independently
encodes the linear conditions
\[
 d_j(a_{u^{-1}(j)}-b_j)=d_0(a_{u^{-1}(0)}-b_0)
\]
at every Boolean vertex. Finite-field ranks at the prime $2147483647$
give first-variation output dimensions $2,4,6$ for $n=1,2,3$, and
$n+1$ for $4\le n\le10$. These ranks are only diagnostic observations,
not characteristic-zero certificates. The preceding proof supplies
the exact real result for every $n\ge4$, including both directions.
No computational observation is used in that proof.

\end{document}
